\subsection{Enhanced recovery and the integral homotopy type}

\begin{theorem}[D006E: \status{Proved}; enhanced shadow recovery]
\label{sa08:D006E}
Let $k$ be a commutative ring, $S=\Gamma_k(T)$, and
$\mathcal G=\operatorname{Gr}_{\geq0}(S)$. Let $\mathcal T$ be the
localizing Serre subcategory generated by modules of bounded weight,
and write $Q:\mathcal G\to\mathcal H=\mathcal G/\mathcal T$ for the
quotient. The functor
\[
 \sigma:D(k)\longrightarrow D(\mathcal H),\qquad
 K\longmapsto Q(K\otimes_k S)
\]
is fully faithful and strong symmetric monoidal on the unbounded
stable derived infinity categories. For every simplicial set $X$ there
is a natural equivalence
\[
 Q C_{\mathrm{gr}}(X)\simeq\sigma C^*(X;k)
 \quad\text{in }\operatorname{CAlg}(D(\mathcal H)).
\]
Here the strict weighted PD-form DGA is regarded as an
$E_\infty$ algebra by restriction along a fixed cofibrant operad
resolution $E_\infty\to\mathrm{Com}$. Its strict algebra structure
supplies this enhancement; no additional operations depending on $X$
are chosen. Consequently the whole $E_\infty$ algebra of ordinary
cochains is recovered from the weighted DGA, uniquely up to a
contractible space of choices.
\end{theorem}

\begin{proof}
We first establish the unbounded categorical assertions. Denote the
section functor of $Q$ by $R$. Exactness of $Q$ implies that $R$
preserves $K$-injective complexes: for an acyclic complex $U$ and a
$K$-injective complex $J$ in $\mathcal H$,
\[
 \operatorname{Hom}^{\bullet}(U,RJ)
 =\operatorname{Hom}^{\bullet}(QU,J)
\]
is acyclic. The termwise identity $QRJ=J$ therefore identifies the
derived counit $Q\mathbf RR$ with the identity. Thus $D(Q)$ is a
localization with kernel
\[
 D_{\mathcal T}(\mathcal G)
 =\{K:H^i(K)\in\mathcal T\text{ for every }i\}.
\]
This is an assertion about mapping complexes and hence about the stable
infinity categories, rather than only their homotopy categories.

The kernel is localizing and stable under upward weight shifts.
Indeed upward shift is exact and preserves colimits, and the inverse
image of $\mathcal T$ under this functor is localizing Serre and
contains every bounded-weight module. Fix $P$ in this kernel. The
full subcategory of $V$ for which $P\otimes_S^{\mathbf L}V$ is in the
kernel is localizing, because derived tensor preserves colimits
separately. It contains $S[-n]$ for every weight $n\geq0$: tensoring
with this object is upward weight shift. These objects, with all
cohomological shifts, compactly generate $D(\mathcal G)$, since
vanishing of their derived Hom groups detects every weightwise
cohomology group. Thus the kernel is a tensor ideal. The symmetric
monoidal localization property endows $D(\mathcal H)$ with the
induced tensor and makes $D(Q)$ strong symmetric monoidal. On modules
this is the tensor induced from graded tensor over $S$; free
resolutions and colimits give its derived extension. In particular no
unbounded tensor spectral sequence or unjustified convergence
assertion is required.

The precise formal statement used here is Lurie, \emph{Higher
Algebra}, Proposition~2.2.1.9, p.~197, with the tensor criterion of
Example~2.2.1.7. These statements were opened and their hypotheses
checked against the preceding tensor-ideal proof.

The exact coefficient-extraction adjunctions proved above are
\[
 \overline L\dashv\sigma\dashv\psi,
 \qquad\psi\sigma=\mathrm{Id}.
\]
If $I$ is $K$-injective over $k$, exactness of $\overline L$ gives
\[
 \operatorname{Hom}^{\bullet}(U,\sigma I)
 =\operatorname{Hom}^{\bullet}(\overline L U,I)
\]
acyclic for every acyclic $U$. Thus $\sigma I$ is $K$-injective.
Applying $\psi\sigma=\mathrm{Id}$ to a $K$-injective resolution
proves the derived unit on all of $D(k)$ and full faithfulness.
Induction $F(K)=K\otimes_k S$ is strong symmetric monoidal:
\[
 F(K)\otimes_S^{\mathbf L}F(L)
 \simeq F(K\otimes_k^{\mathbf L}L).
\]
The $k$-flatness of $S$ identifies this induction with the termwise
functor, and composition with $Q$ proves the assertion for $\sigma$.
A fully faithful symmetric monoidal functor is fully faithful on
commutative algebra objects: all tensor powers remain in its image,
and every structure map and its coherence data lift through the
corresponding equivalences of mapping spaces.
This is also Lurie, \emph{Higher Algebra}, Remark~2.2.1.5,
p.~196, applied to the full monoidal essential image.

We construct the comparison as an $E_\infty$ map before proving that
its underlying map becomes an equivalence. Let $I_X=\Delta/X$ be
the category of simplices, with projection $p:I_X\to\Delta$.
Consider the diagram of strict DGAs
\[
 I_X^{\mathrm{op}}\longrightarrow\mathrm{CDGA}(\mathcal G),
 \qquad (\Delta[n]\xrightarrow{x}X)\longmapsto\mathcal A_n^\bullet.
\]
Its ordinary limit is $C_{\mathrm{gr}}(X)$: limits in $\mathcal G$
are weightwise. Evaluation at every simplex is a strict DGA map and
forms a compatible cone. Restricting along $E_\infty\to\mathrm{Com}$
and passing to derived algebras gives a cone of $E_\infty$ algebras,
hence a canonical map
\begin{equation}\label{eq:D-einfinity-cone}
 C_{\mathrm{gr}}(X)
 \longrightarrow\operatorname*{holim}_{I_X^{\mathrm{op}}}
       \mathcal A_n^\bullet.
\end{equation}
This uses an $E_\infty$ algebra model, and does not require a model
structure on strict commutative DGAs in positive characteristic.

The PD Poincare contraction shows, in every simplicial degree, that
$S\to\mathcal A_n^\bullet$ is a quasi-isomorphism of strict DGAs.
It is natural in simplicial operators, so the diagram on the right
of \eqref{eq:D-einfinity-cone} is equivalent to the constant diagram
$S$. The nerve of $I_X$ has the homotopy type of $X$. The homotopy
limit of the constant $E_\infty$ algebra $k$ over this space is its
ordinary $E_\infty$ cochain algebra $C^*(X;k)$. To commute $F$ with
this limit, use its exact left adjoint $L$: if $I$ is $K$-injective,
then $FI$ is $K$-injective by
\[
 \operatorname{Hom}^{\bullet}(U,FI)
 =\operatorname{Hom}^{\bullet}(LU,I).
\]
Thus $L\dashv F$ derives on unbounded categories, and the derived
$F$ is a right adjoint and preserves all homotopy limits. Its strong
symmetric monoidality then implies preservation of limits of
commutative algebras, whose underlying limits are those of the
ambient category. We obtain
\[
 \operatorname*{holim}_{I_X^{\mathrm{op}}}\mathcal A_n^\bullet
 \simeq \operatorname*{holim}_{I_X^{\mathrm{op}}}S
 \simeq F C^*(X;k)
\]
as $E_\infty$ algebras.
For the assertion that limits of algebra objects are computed on
their underlying objects, the opened locators are Lurie,
\emph{Higher Algebra}, Corollaries~3.2.2.4--3.2.2.5, pp.~358--359.

The underlying map in \eqref{eq:D-einfinity-cone} is precisely the
ordinary-to-derived sections comparison. One can check this without
identifying two unspecified resolutions: restriction $p^*$ on module
presheaves preserves injectives because its left adjoint is, at
$[n]$, the direct sum of the modules indexed by $x\in X_n$, and is
exact. Hence $\mathbf R\lim p^*$ computes $\mathbf R\Gamma_X$.
In weight $w$, the dimension-shifting comparison is an isomorphism
through cohomological degree $w$; the form complex has no terms in
higher degrees. The cone of the underlying map therefore has
bounded-weight cohomology in every degree. It lies in
$D_{\mathcal T}(\mathcal G)$. Applying the strong symmetric monoidal
localization $Q$ to \eqref{eq:D-einfinity-cone} gives the asserted
$E_\infty$ equivalence. Full faithfulness on commutative algebras
completes the recovery claim.
\end{proof}

\paragraph{Strict commutativity and the recovered higher operations.}
The monoidal category in this theorem is the category of weight-graded
$S$-modules, followed by its derived shadow localization. The functor
forgetting weights does not factor through that localization: a nonzero
module of bounded weight is killed by $Q$ but remains nonzero after
forgetting weights. Likewise, exact coefficient extraction is not a
strong symmetric monoidal functor. Its colax tensor comparison on two
weight-one free generators is
\[
 L(S[-2])\longrightarrow L(S[-1])\otimes_k L(S[-1]),
 \qquad k\xrightarrow{\ 2\ }k.
\]
Indeed the two adjunction units send the generators to $e_1$, and
$e_1e_1=2e_2$; the normalized generator of $L(S[-2])$ maps to one.
In characteristic two this comparison is zero. Applying $L$ termwise
therefore does not produce a strict commutative $k$-DGA with the
recovered operations. For example,
$H^1(C_{\mathrm{gr}}(S^1))=I_1$ has no weight-zero part, whereas
$QI_1\simeq QS$ because $S/I_1$ has bounded weight. Its recovered
degree-one class uses that localization equivalence. Higher operations
on such classes are determined by the coherent derived comparison
above; they are not computed by forgetting weights or by transporting
only products of homogeneous representatives.

\begin{corollary}[D006M: \status{Known, with the proved comparison};
 integral homotopy determination]\label{sa08:D006M}
For nilpotent spaces of finite type, the strict weighted integral
PD-form DGA determines the integral weak homotopy type. More
precisely, if the corresponding weighted DGAs become equivalent as
commutative algebra objects after the indicated derived shadow
localization, then the spaces are weakly equivalent. This does not
assert fullness of the functor on morphisms.
\end{corollary}
\begin{proof}
Theorem~\ref{sa08:D006E} identifies their recovered algebras with
$C^*(X;\mathbb Z)$ and $C^*(Y;\mathbb Z)$ as $E_\infty$ algebras.
Mandell's \emph{Cochains and homotopy type}, Main Theorem, says that
two finite-type nilpotent spaces are weakly equivalent if and only
if their integral cochain algebras are quasi-isomorphic as
$E_\infty$ algebras. Apply that theorem to the recovered equivalence.
Conversely a weak equivalence of spaces induces an equivalence of
the recovered cochain algebras and hence of their shadow algebras.
\end{proof}

\paragraph{Scope.}
This uses the whole strict weighted cochain DGA and its canonical
commutative-algebra enhancement. It does not infer higher operations
from a binary cohomology product, or from an algebra object in the
ordinary homotopy category alone. It applies to the graded convention
$\bigoplus_w\operatorname{Hom}(X,\mathcal A_w)$; it does not replace
that object by literal unbounded-weight cochains on infinite $X$.
The separate proposed binomial associated-graded comparison is not
needed for this recovery theorem.

\paragraph{Opened comparisons and locators.}
The nerve-of-simplices fact is Hirschhorn, \emph{Model Categories and
Their Localizations}, Theorem~18.9.3; the author's opened paper
\emph{The Diagonal of a Multicosimplicial Object}, proof of
Proposition~9.6, p.~20, explicitly states and uses that theorem.
The full book itself was not opened. Richter--Sagave,
\emph{A strictly commutative model for the cochain algebra of a
space}, Theorems~1.1, 1.2 and 6.2, provide a related construction
using commutative $\mathcal I$-diagrams. Their result was opened and
checked as a comparison, and is not cited as a proof of the
particular shadow theorem established here. No priority is claimed.
